\section{Comprensión del problema}

\subsection{Problema abordado}
Se busca analizar características clínicas medidas en tumores mamarios para diferenciar entre diagnósticos \textbf{benignos (B)} y \textbf{malignos (M)}, apoyando una clasificación temprana y confiable.

\subsection{Decisión que se puede apoyar}
Con este análisis se pueden priorizar casos con mayor probabilidad de malignidad para evaluación médica más rápida, apoyo al tamizaje y seguimiento clínico.

\subsection{Usuarios del análisis}
Este análisis puede ser usado por personal médico (oncología, radiología, medicina general), equipos de apoyo diagnóstico, investigadores biomédicos y gestores hospitalarios.

\subsection{Importancia}
La detección oportuna impacta directamente el pronóstico y tratamiento del cáncer de mama. Un análisis adecuado puede reducir retrasos y mejorar el uso de recursos clínicos.

\subsection{Preguntas guía}
\begin{itemize}
    \item \textbf{Tipo de problema:} principalmente predictivo, con apoyo descriptivo en la etapa exploratoria.
    \item \textbf{Variable más importante:} \texttt{diagnosis}, al ser la variable objetivo.
    \item \textbf{Riesgo de mala interpretación:} decisiones clínicas erradas, especialmente falsos negativos (clasificar un caso maligno como benigno).
\end{itemize}
